% 03-sun-improved — recap of the Hamiltonian + the van de Sande subtraction
\section{SU(N) in 1+1 dimensions and the improved Hamiltonian}
\label{sec:sun}

The theory is SU($N$) gauge theory coupled to a single fundamental fermion of
mass $m$,
\begin{equation}
  \mathcal{L} \;=\; -\tfrac14 F^a_{\mu\nu}F^{a\,\mu\nu}
  \;+\; \bar\psi\,(i\slashed{D} - m)\,\psi ,
  \label{eq:lagrangian}
\end{equation}
in one space and one time dimension, where the coupling $g$ carries dimensions
of mass and $m/g$ is the only physical parameter.  In the gauge $A^+ = 0$ the
gluon field carries no dynamical degrees of freedom: $A^-$ is fixed by a
constraint equation, and integrating it out leaves quarks interacting through
an instantaneous linear potential in $x^-$ --- confinement is kinematic in two
dimensions, in QED$_{1+1}$ as much as here.  Half of the fermion components
are likewise constrained, so the physical content is a single spinless quark
field per colour.

The light-cone Hamiltonian then takes the form
\begin{equation}
  P^- \;=\; m^2 \sum_{n,c} \frac{a^\dagger_{n c} a_{n c}}{k^+_n}
        \;+\; g^2\, H_I ,
  \label{eq:pminus}
\end{equation}
where $H_I$ contains the normal-ordered four-quark vertices generated by the
instantaneous potential, with the colour structure
$\delta_{c_1 c_3}\delta_{c_2 c_4} - \tfrac1N \delta_{c_1 c_2}\delta_{c_3 c_4}$
and a $1/(\Delta k^+)^2$ momentum kernel, together with a one-body
``self-induced inertia'' --- the normal-ordering remnant of the instantaneous
exchange, diagonal in the Fock basis.  Explicit matrix elements are given in
Ref.~\cite{Hornbostel:1988fb}; our implementations of them, in both the
preserved 1988-era Fortran and an independently validated Python port, are
public (Appendix~\ref{app:data}).

\subsection{The endpoint problem and the improved Hamiltonian}
\label{sec:improved}

Near the endpoints $x \to 0, 1$ the continuum wavefunction of a bound state
behaves as $\phi \sim x^{a}$, with the exponent set by the transcendental
condition~\cite{tHooft:1974hx,Hornbostel:1988fb}
\begin{equation}
  1 - \pi a \cot(\pi a) \;=\; \frac{\pi m^2}{g^2 C_F},
  \qquad C_F = \frac{N^2-1}{2N} .
  \label{eq:endpoint-exponent}
\end{equation}
As $m/g \to 0$ the exponent $a \to 0$ and the endpoint region
$x \lesssim e^{-1/a}$ shrinks below any fixed grid: a uniform-in-$k^+$
discretization whose smallest momentum fraction is $1/K$ then replaces the
endpoint integral $\int dx\, x^{2a-1} \sim 1/2a$ by a logarithm $\ln K$,
and no reachable resolution recovers it (Sec.~\ref{sec:budget}).  This is a
property of the grid, not of light-cone quantization: a continuum 't~Hooft
solver on a uniform mesh, sharing no code with the discretized theory, fails
in exactly the same way.

Van de Sande's improved Hamiltonian~\cite{vandeSande:1996mc} treats the
disease at its source.  The self-induced inertia and the Coulomb exchange
suffer separate, cancelling endpoint singularities; the improvement replaces
the free-field self-energy by one computed against the interacting endpoint
behaviour of Eq.~\eqref{eq:endpoint-exponent}, so that the residual kernel
vanishes when either momentum approaches an endpoint.  Two properties make
this essentially free within an existing DLCQ code.  First, the modification
is exactly diagonal: it amounts to swapping the self-energy table
$\sigma_{\rm std} \to \sigma_{\rm imp}$, leaving every interaction kernel ---
and the standard path, bit for bit --- untouched.  Second, the colour weight
it requires is $\langle -T_a\!\cdot\!T_c\rangle$ at zero momentum transfer,
which the colour-contraction machinery already computes; the identity
$\sum_{c\neq a}\langle -T_a\!\cdot\!T_c\rangle = C_F$ holds exactly in our
implementation at every $N$ and parton number we use, and is enforced as a
gate rather than assumed.

One rule attaches to everything that follows and is worth stating once,
prominently.  The two Hamiltonians converge to the continuum along different
paths: the improved spectrum is extrapolated in a plain $1/K$ series
\cite{vandeSande:1996mc}, while the standard spectrum requires the
fractional-power ladder of Ref.~\cite{Hornbostel:1988fb}, Eq.~(27) there.
The two bases are never mixed, and no fit in this paper crosses them; the
error budget in Sec.~\ref{sec:budget} quantifies what happens when one does.
